\begin{frame}[label=cierre]{Qué preguntas respondimos y cómo se sostienen}
  \small
  \begin{tabularx}{\linewidth}{@{}lXX@{}}
    \toprule
    \textbf{Ítem} & \textbf{Pregunta resuelta} & \textbf{Respuesta verificable}\\
    \midrule
    3b & ¿Cuántas personas por fila? & Promedio 1.2941, mediana 1.\\
    3c.1 & ¿Qué peso tienen los menores? & 57.79\% nacional, 72.21\% en Madre de Dios.\\
    3c.2 & ¿Qué peso tienen los extranjeros? & 2.72\% en 2024.\\
    3c.3 & ¿Qué meses superan la referencia uniforme? & Diciembre: 9.49\%, índice 1.14.\\
    5e.2 & ¿Qué perfil concentra más personas reportadas? & Adolescentes mujeres: 51\,021 personas, 42.02\%.\\
    4 y 5g & ¿Cómo conservar y entregar los resultados? & Base y dump: 121\,427; seis libros Excel registrados.\\
    \bottomrule
  \end{tabularx}
  \begin{block}{La evidencia que conecta pregunta y respuesta}
    \footnotesize Cada ficha muestra el ítem, las columnas, las operaciones, el código y su salida.
    Las 25 consultas completas están disponibles desde el \enlace{mapa}{mapa de consultas}.
  \end{block}
  \footnotesize
  Las conclusiones describen lo reportado en el dataset. Las proporciones no estiman
  riesgo poblacional, y los cambios temporales no demuestran causalidad.
  \note{40 s. Resumir las respuestas y cerrar aquí el recorrido oral. Total sugerido: 570 segundos, con 30 segundos de margen. Los anexos solo se abren por consulta.}
\end{frame}
